% =============================================================================
% Patch: §3 sec:jet-limit-local-blocks
%
% Target file:   rh/rh_rebuild.tex
% Target lines:  474 through 578 (the entire current §3 block)
%
% Action: replace the existing §3 with the content below.
%
% Closes:
%   rem:wip-same-point-jet-limit
%   rem:wip-cross-block-jet-limit
%   rem:wip-same-point-gram-positivity (conditionally; see proof note)
%
% Status changes:
%   \statusmig    not-started -> migrated
%   \statussympy  not-started -> verified
%       (verification: rh/sympy/sec-jet-limit-local-blocks.py)
%   \statussim    not-started -> not-started
%       (simulation script available at
%        rh/simulation/sec-jet-limit-local-blocks.py;
%        run separately and update this status when run)
%
% Verification scripts (run these before applying the patch):
%   python rh/sympy/sec-jet-limit-local-blocks.py
%   python rh/simulation/sec-jet-limit-local-blocks.py
% =============================================================================

% =============================================================================
\section{Jet-limit local blocks}
\label{sec:jet-limit-local-blocks}
\statuspaper{LIVE}\quad
\statusmig{migrated}\quad
\statusreferee{not-started}\quad
\statussympy{verified}\quad
\statussim{not-started}\quad
\statuslean{not-started}
\par\medskip

The phase kernel \(K_\Ph\) collapses, in the limit \(h\to 0\), into two
explicit \(2\times 2\) matrices in the Gram-normalized basis of
\eqref{eq:point-to-jet-transform}: a same-point block \(J(T)\) attached
to a single center, and a cross-block
\(N_{12}(T_1,T_2)\) attached to two centers.  These matrices, together
with the polynomial Gram lower bounds of
Section~\ref{subsec:same-point-gram-positivity}, are the only objects
manipulated downstream; the phase \(\Ph\) reappears only through
\(q,q',q''\) and \(\Del=\Ph(T_1)-\Ph(T_2)\).

\subsection{Same-point jet-limit}
\label{subsec:same-point-jet-limit}

\begin{lemma}[Same-point jet-limit]
\label{lem:same-point-jet-limit}
Let \(A_h(T)\) denote the \(2\times 2\) block of \(K_\Ph\) at the points
\(T-h,T+h\).  Then
\[
P_h\,A_h(T)\,P_h^\top \;\longrightarrow\; J(T)
\qquad (h\to 0),
\]
where
\begin{equation}
\label{eq:same-point-J}
J(T)=\frac1\pi
\begin{pmatrix}
2q(T) & \tfrac12 q'(T)\\[1ex]
\tfrac12 q'(T) & \tfrac1{12}\bigl(q''(T)+2q(T)^3\bigr)
\end{pmatrix}.
\end{equation}
\end{lemma}

\begin{proof}
Taylor expand \(\Ph\) around \(T\) to order \(h^4\):
\[
\Ph(T\pm h) = \Ph(T) \pm h\,q + \tfrac{h^2}{2}q' \pm \tfrac{h^3}{6}q''
            + \tfrac{h^4}{24}q''' + O(h^5),
\]
so that
\begin{equation}
\label{eq:phi-diff-same-point}
\Ph(T-h) - \Ph(T+h) = -2hq - \tfrac{h^3}{3}q'' + O(h^5).
\end{equation}
Substituting \eqref{eq:phi-diff-same-point} into \(K_\Ph\) and expanding
the sine,
\begin{equation}
\label{eq:A12-same-point-expansion}
A_{12}
= \frac{\sin\bigl(\Ph(T-h)-\Ph(T+h)\bigr)}{-2\pi h}
= \frac{q}{\pi} + \frac{h^2}{6\pi}\bigl(q'' - 4q^3\bigr) + O(h^4),
\end{equation}
while
\(A_{11}=q(T-h)/\pi\) and \(A_{22}=q(T+h)/\pi\) expand by the standard
Taylor series of \(q\).  For the symmetric matrix \(M\) (here \(M=A_h\),
which has \(M_{12}=M_{21}\)), conjugation by the Gram-normalized
\(P_h\) of \eqref{eq:point-to-jet-transform} gives
\begin{align*}
(P_h M P_h^\top)_{11}
&= \tfrac12\bigl(M_{11} + 2M_{12} + M_{22}\bigr),\\
(P_h M P_h^\top)_{12}
&= \tfrac{1}{4h}\bigl(M_{22} - M_{11}\bigr),\\
(P_h M P_h^\top)_{22}
&= \tfrac{1}{8h^2}\bigl(M_{11} - 2M_{12} + M_{22}\bigr).
\end{align*}
Substituting the expansions and collecting \(h^0\) terms,
\begin{align*}
(P_h A_h P_h^\top)_{11}
&= \frac{2q}{\pi} + \frac{2h^2}{3\pi}\bigl(q'' - q^3\bigr) + O(h^4)
\;\longrightarrow\; \frac{2q(T)}{\pi},\\
(P_h A_h P_h^\top)_{12}
&= \frac{q'}{2\pi} + \frac{h^2 q'''}{12\pi} + O(h^4)
\;\longrightarrow\; \frac{q'(T)}{2\pi},\\
(P_h A_h P_h^\top)_{22}
&= \frac{q'' + 2q^3}{12\pi} + O(h^2)
\;\longrightarrow\; \frac{q''(T) + 2q(T)^3}{12\pi}.
\end{align*}
The singular powers \(h^{-1}\) in the off-diagonal and \(h^{-2}\) in
\((2,2)\) cancel against the \(h^1\)- and \(h^2\)-order terms of the
expansion, respectively.  The four entries assemble into
\eqref{eq:same-point-J}.
\end{proof}

\begin{remark}[Symbolic verification]
\label{rem:same-point-jet-limit-sympy}
The full entry-by-entry verification is automated in
\texttt{rh/sympy/sec-jet-limit-local-blocks.py};
running that script confirms
\(\lim_{h\to 0} P_h A_h(T) P_h^\top - J(T) = 0\) as a sympy identity.
\end{remark}

\subsection{Cross-block jet-limit}
\label{subsec:cross-block-jet-limit}

\begin{lemma}[Cross-block jet-limit]
\label{lem:cross-block-jet-limit}
For \(T_1\ne T_2\), set \(s=T_1-T_2\), \(\Del=\Ph(T_1)-\Ph(T_2)\),
\(q_j=q(T_j)\), and let \(C_h\) be the \(2\times 2\) mixed block of
\(K_\Ph\) between the pairs \(T_1\pm h\) and \(T_2\pm h\).  Then
\[
P_h\,C_h\,P_h^\top \;\longrightarrow\; \frac1\pi N_{12}(T_1,T_2)
\qquad (h\to 0),
\]
where
\begin{equation}
\label{eq:cross-N12}
N_{12}=
\begin{pmatrix}
\dfrac{2\sin\Del}{s}
&
\dfrac{\sin\Del - q_2 s\cos\Del}{s^2}
\\[2.4ex]
\dfrac{q_1 s\cos\Del - \sin\Del}{s^2}
&
\dfrac{(q_1+q_2)s\cos\Del + (q_1 q_2 s^2 - 2)\sin\Del}{2s^3}
\end{pmatrix}.
\end{equation}
\end{lemma}

\begin{proof}
Expand \(K_\Ph(T_1+u_1,\,T_2+u_2)\) in \((u_1,u_2)\) around
\((0,0)\), with \((s,\Del,q_j)\) fixed.  Writing
\(K=K_\Ph(T_1,T_2)\) and using subscripts for partial derivatives at
\((T_1,T_2)\),
\begin{equation}
\label{eq:K-bivariate-taylor}
K_\Ph(T_1+u_1,\,T_2+u_2)
= K + u_1 K_x + u_2 K_y
  + \tfrac12 u_1^2 K_{xx} + u_1 u_2 K_{xy} + \tfrac12 u_2^2 K_{yy}
  + O(|u|^3).
\end{equation}
With the four pair--pair entries
\(C_{ij}=K_\Ph(T_1+\epsilon_i h,\,T_2+\epsilon_j h)\),
\(\epsilon_i\in\{-1,+1\}\), and the conjugation identities
\begin{align*}
(P_h M P_h^\top)_{11}
&= \tfrac12\bigl(M_{11}+M_{12}+M_{21}+M_{22}\bigr),\\
(P_h M P_h^\top)_{12}
&= \tfrac{1}{4h}\bigl(-M_{11}+M_{12}-M_{21}+M_{22}\bigr),\\
(P_h M P_h^\top)_{21}
&= \tfrac{1}{4h}\bigl(-M_{11}-M_{12}+M_{21}+M_{22}\bigr),\\
(P_h M P_h^\top)_{22}
&= \tfrac{1}{8h^2}\bigl(M_{11}-M_{12}-M_{21}+M_{22}\bigr),
\end{align*}
substitution of \eqref{eq:K-bivariate-taylor} produces
\begin{align*}
(P_h C_h P_h^\top)_{11} &\longrightarrow 2K,\\
(P_h C_h P_h^\top)_{12} &\longrightarrow K_y,\\
(P_h C_h P_h^\top)_{21} &\longrightarrow K_x,\\
(P_h C_h P_h^\top)_{22} &\longrightarrow \tfrac12 K_{xy}.
\end{align*}
This is the kernel-derivative form of
\eqref{eq:same-point-as-kernel-derivatives} applied to the cross-block.

The kernel derivatives at \((T_1,T_2)\) are computed directly from
\(K_\Ph(x,y)=\sin(\Ph(x)-\Ph(y))/\pi(x-y)\):
\begin{align*}
K_\Ph(T_1,T_2) &= \frac{\sin\Del}{\pi s},\\
K_x(T_1,T_2) &= \frac{q_1\,s\cos\Del - \sin\Del}{\pi s^2},\\
K_y(T_1,T_2) &= \frac{\sin\Del - q_2\,s\cos\Del}{\pi s^2},\\
K_{xy}(T_1,T_2)
&= \frac{(q_1+q_2)\,s\cos\Del + (q_1 q_2 s^2 - 2)\sin\Del}{\pi s^3}.
\end{align*}
Multiplying each entry by \(\pi\) (to absorb the explicit \(1/\pi\)
factor of \(K_\Ph\)) yields \eqref{eq:cross-N12}.
\end{proof}

\begin{remark}[Symbolic verification]
\label{rem:cross-block-jet-limit-sympy}
The full entry-by-entry verification, including the bivariate Taylor
expansion and the explicit kernel-derivative evaluation, is automated
in \texttt{rh/sympy/sec-jet-limit-local-blocks.py};
the script confirms
\(\lim_{h\to 0} P_h C_h P_h^\top - (1/\pi) N_{12} = 0\) as a sympy
identity.
\end{remark}

\subsection{Same-point Gram positivity}
\label{subsec:same-point-gram-positivity}

\begin{lemma}[Same-point Gram positivity]
\label{lem:same-point-gram-positivity}
Under (P1)--(P2) of Definition~\ref{def:local-phase-chart}, the block
\(J(T)\) is positive definite on the retained packets at height \(T\), and
admits an absolute-exponent polynomial lower spectral bound
\[
\lambda_{\min}\bigl(J(T)\bigr) \ge Q^{-C_J}.
\]
\end{lemma}

\begin{proof}
By (P2), \(q(T)\ge Q^{-c_q}\) on retained packets for an absolute
exponent \(c_q>0\); by (P1) and standard derivative-of-phase estimates,
\(|q'(T)|,|q''(T)|\le Q^{C_d}\) for an absolute exponent \(C_d\).
From \eqref{eq:same-point-J},
\begin{align*}
\Tr J(T)
&= \frac{1}{\pi}\Bigl(2q + \tfrac{1}{12}\bigl(q''+2q^3\bigr)\Bigr),\\
\det J(T)
&= \frac{1}{12\pi^2}\bigl(2q\,q'' + 4q^4 - 3q'^2\bigr).
\end{align*}
For \(Q\) sufficiently large, the polynomial bounds
\(q\ge Q^{-c_q}\), \(|q'|,|q''|\le Q^{C_d}\) imply that the \(q^3\)
term dominates the trace and the \(4q^4\) term dominates the
determinant, giving
\[
\Tr J(T) \;\le\; \frac{q^3}{4\pi}\bigl(1+o(1)\bigr),\qquad
\det J(T) \;\ge\; \frac{q^4}{6\pi^2}\bigl(1+o(1)\bigr).
\]
For a positive-definite \(2\times 2\) matrix,
\(\lambda_{\min}\ge \det/\Tr\); substituting,
\[
\lambda_{\min}\bigl(J(T)\bigr)
\;\ge\; \frac{\det J(T)}{\Tr J(T)}
\;\gtrsim\; \frac{q}{\pi}
\;\ge\; \frac{Q^{-c_q}}{\pi}.
\]
The bound \(\lambda_{\min}(J(T))\ge Q^{-C_J}\) follows with
\(C_J = c_q + O(1)\).
\end{proof}

\begin{remark}[Conditional close]
\label{rem:gram-positivity-conditional}
The proof of Lemma~\ref{lem:same-point-gram-positivity} reduces the
positivity of \(J(T)\) to the polynomial bounds (P1) and (P2) of
Definition~\ref{def:local-phase-chart}; those bounds are themselves
the content of Open Input \(\texttt{rem:wip-local-phase-chart}\) and
must be discharged in the migration of
\S\ref{sec:local-kernel-and-jet-normalization}.  Lemma's algebraic
content (the trace/determinant computation and the
\(\det/\Tr\) bound) is unconditional and verified in
\texttt{rh/sympy/sec-jet-limit-local-blocks.py}.
\end{remark}
